\chapter{Actualizaciones}

Un release tiene partes que se actualizan por caminos distintos:

\begin{center}
\begin{tabularx}{\linewidth}{@{}l l X@{}}
\toprule
\textbf{Parte} & \textbf{Dónde vive} & \textbf{Cómo se actualiza} \\
\midrule
Bitstream de la FPGA & flash de configuración de la FPGA & JTAG (FlashPro), con acceso físico \\
Driver y servicio & tarjeta microSD & por red, con \cmd{mpeg2fpga-update} (paquete firmado) \\
Librería cliente & PC del usuario & \cmd{pip install} del nuevo wheel \\
\bottomrule
\end{tabularx}
\end{center}

\begin{advertencia}[Limitación de la placa]
El PolarFire SoC permite actualizar el bitstream en campo (\emph{Auto Update}
e \emph{IAP}), pero ambos requieren una memoria flash SPI conectada al System
Controller del chip, que la Discovery Kit no tiene. En esta placa el
bitstream sólo se actualiza por JTAG. Una placa de producto debería incluir
esa memoria.
\end{advertencia}

\section{Compatibilidad}

\begin{itemize}
  \item \cmd{/v1/device} informa las versiones de las tres partes de la placa;
    \cmd{bitstream} identifica el release y el commit exacto.
  \item Dentro del protocolo v1, un cliente de cualquier versión funciona con
    un dispositivo de cualquier versión: el protocolo sólo crece por adición.
\end{itemize}

\section{Actualizar el software de la placa}

Con el paquete y su firma copiados a la placa, como root:

\begin{lstlisting}[language=shell]
scp mpeg2fpga-board-<versión>.tar.gz* root@placa:/tmp/
ssh root@placa /usr/local/sbin/mpeg2fpga-update /tmp/mpeg2fpga-board-<versión>.tar.gz
\end{lstlisting}

\cmd{mpeg2fpga-update} verifica la firma del paquete contra la clave del
proyecto que quedó en la placa en la primera instalación
(\cmd{/etc/mpeg2fpga/release-key.pub.pem}), nunca contra una que venga en el
paquete. Si la firma falta o no corresponde, no instala nada. Si corresponde,
recarga el driver y reinstala el servicio, conservando sus opciones
(\cmd{-{}-tls}, \cmd{-{}-debug}), el token y el certificado.

\section{Verificar un release}
\label{sec:verificar}

Cada release publica \cmd{SHA256SUMS}, con la suma de cada archivo, y
\cmd{SHA256SUMS.sig}, su firma Ed25519 con la clave del proyecto. La clave
pública, \cmd{mpeg2fpga-release.pub.pem}, se toma del repositorio de fuentes
(\cmd{release/keys/}), no del mismo lugar que los archivos. Con OpenSSL~3:

\begin{lstlisting}[language=shell]
openssl pkeyutl -verify -pubin -inkey mpeg2fpga-release.pub.pem \
    -rawin -in SHA256SUMS -sigfile SHA256SUMS.sig
sha256sum -c SHA256SUMS
\end{lstlisting}

La primera orden prueba que las sumas son del proyecto; la segunda, que cada
archivo coincide con ellas.
